O desbalanceamento de classes é frequentemente tratado como o principal obstáculo à predição de desfechos raros, embora acrescentar observações minoritárias não recupere uma fronteira de decisão fraca. Esta qualificação investiga se a separabilidade de classes deve ser um objetivo explícito da geração de dados sintéticos tabulares de saúde. Experimentos iniciais de reamostragem, aumento com CTGAN e filtragem por âncoras produziram ganhos limitados ou inconsistentes, motivando uma análise da sobreposição entre classes. O método SepAware gera um conjunto ampliado de candidatos CTGAN e seleciona dados sintéticos, preservando as contagens de classe, segundo realismo, separabilidade e compatibilidade com âncoras de domínio. Todos os resultados preditivos usam validação cruzada estratificada de cinco partes com controle de vazamento, em dados de obesidade do Vigitel e em 334 registros hospitalares de COVID-19. No experimento fatorial $2^2$, separabilidade explica 96,00\% da variação entre tratamentos e erro no Macro-F1 treino-sintético/teste-real do Vigitel ($p=5,03\times10^{-10}$), mas apenas 5,13\% em COVID-19, onde o erro explica 81,53\%. Os resultados sustentam um mecanismo dependente do conjunto de dados: separabilidade explícita pode recuperar sinal condicional útil, mas também pode selecionar protótipos fáceis e eliminar sobreposição clinicamente importante. O trabalho restante testará robustez entre sementes, classificadores, geradores, níveis de desbalanceamento e bases externas, juntamente com diversidade minoritária, consistência causal, plausibilidade clínica e privacidade.
